\section{账本事件的社会制度深描：chenghua-zhengde}

\label{sec:ledger-depth-chenghua-zhengde}

以下各节依本卷账本的事件、来源和置信提示组织，不以编年节点替代地方社会。朝廷命令、法律条文、城市记录、家庭文书与物质遗存的证明范围不同，必须让其差异保留在叙事中。

\subsection{政治制度与地方执行}

在账本条目\texttt{undefined}中，1464（天顺八年）正月的“英宗去世，皇太子朱见深即位，是为宪宗”发生于明。本条把背景说明为朱见深已于天顺二年复立为太子，继承程序较景泰末年稳定。，并以成化朝得以在既有储位基础上开局，但宦官、后宫与边疆军务将构成新的权力场域。概括可见影响。要理解它的社会后果，不能止于宫廷或战场：土地、赋役、司法、道路、性别化家庭劳动、城市市场、书院寺观和边海网络都会影响地方如何接收或改写制度。

本条依据《英宗实录》天顺八年正月条；《宪宗实录》成化元年即位条；《明史·宪宗本纪》。这些来源能够较稳妥地支持所记年序、诏令、任命或特定地点的行动，但无法自动证明全国的财产、人口、识字、信仰或动机。账本保留的证据说明是“死亡、即位与年号转换可靠；先朝遗诏和临终安排的细节须注意实录的礼制化记录。”。所以，官府标签、奏疏数字与后来的道德评语必须放回其文书目的，不能被写成不经分区核验的社会全貌。

这条事件更适合作为一组关系变化的锚点：中央方案需由地方官、书吏、士绅、宗族、商人与劳动者共同转译；不同家庭面对的税役、婚姻、迁徙和安全风险并不相等。材料只让我们看见其中有限的记录者。承认可见与不可见的界线，才能使制度史、文化史和区域史互相校正。

在账本条目\texttt{undefined}中，1465（成化元年）的“蒙古诸部、朝鲜及西北绿洲的交涉在本年留下多方材料，应以具体使团和地名校正概称。”发生于明。本条把背景说明为朝贡名义、边境安全与港口贸易的利益使交涉持续发生。，并以它为后续册封、互市或海防安排留下文书与跨政权比较线索。概括可见影响。要理解它的社会后果，不能止于宫廷或战场：土地、赋役、司法、道路、性别化家庭劳动、城市市场、书院寺观和边海网络都会影响地方如何接收或改写制度。

本条依据《宪宗实录》成化元年条；《明史·外国传》；《朝鲜王朝实录》、琉球《历代宝案》或海外档案。这些来源能够较稳妥地支持所记年序、诏令、任命或特定地点的行动，但无法自动证明全国的财产、人口、识字、信仰或动机。账本保留的证据说明是“桥接条目只标示可回查的年度问题与材料链；实录有中央选择性，崇祯期尤其需要奏疏、地方、朝鲜及满文材料交叉。”。所以，官府标签、奏疏数字与后来的道德评语必须放回其文书目的，不能被写成不经分区核验的社会全貌。

这条事件更适合作为一组关系变化的锚点：中央方案需由地方官、书吏、士绅、宗族、商人与劳动者共同转译；不同家庭面对的税役、婚姻、迁徙和安全风险并不相等。材料只让我们看见其中有限的记录者。承认可见与不可见的界线，才能使制度史、文化史和区域史互相校正。

在账本条目\texttt{undefined}中，1465（成化元年）的“地方灾异、赈济与水利条目为本年社会史提供入口，地方志的修纂层次需要说明。”发生于明。本条把背景说明为自然风险与地方报告促成朝廷或地方的应对记录。，并以赈济、迁徙和损失的实际范围仍须以受灾地区材料分别核验。概括可见影响。要理解它的社会后果，不能止于宫廷或战场：土地、赋役、司法、道路、性别化家庭劳动、城市市场、书院寺观和边海网络都会影响地方如何接收或改写制度。

本条依据《宪宗实录》成化元年条；《明史·五行志》；受灾府县地方志与奏疏。这些来源能够较稳妥地支持所记年序、诏令、任命或特定地点的行动，但无法自动证明全国的财产、人口、识字、信仰或动机。账本保留的证据说明是“桥接条目只标示可回查的年度问题与材料链；实录有中央选择性，崇祯期尤其需要奏疏、地方、朝鲜及满文材料交叉。”。所以，官府标签、奏疏数字与后来的道德评语必须放回其文书目的，不能被写成不经分区核验的社会全貌。

这条事件更适合作为一组关系变化的锚点：中央方案需由地方官、书吏、士绅、宗族、商人与劳动者共同转译；不同家庭面对的税役、婚姻、迁徙和安全风险并不相等。材料只让我们看见其中有限的记录者。承认可见与不可见的界线，才能使制度史、文化史和区域史互相校正。

在账本条目\texttt{undefined}中，1465（成化元年）的“韩雍率军进攻大藤峡，镇压当地反抗并重整广西军政”发生于明与广西大藤峡地区。本条把背景说明为广西山地社群与卫所、土司及州县长期存在赋役和控制冲突，朝廷将其界定为需要军事处置的“叛乱”。，并以明廷增设和调整广西防务、行政节点；山地社会的土地、迁徙和征役后果仍需地方材料追踪。概括可见影响。要理解它的社会后果，不能止于宫廷或战场：土地、赋役、司法、道路、性别化家庭劳动、城市市场、书院寺观和边海网络都会影响地方如何接收或改写制度。

本条依据《宪宗实录》成化元年广西条；《明史·韩雍传》；广西地方志与大藤峡相关碑刻。这些来源能够较稳妥地支持所记年序、诏令、任命或特定地点的行动，但无法自动证明全国的财产、人口、识字、信仰或动机。账本保留的证据说明是“征讨、主帅和朝廷建置可由实录与传记互证；官文中的族名、战果和“平定”不等于当地社群已被完全控制。”。所以，官府标签、奏疏数字与后来的道德评语必须放回其文书目的，不能被写成不经分区核验的社会全貌。

这条事件更适合作为一组关系变化的锚点：中央方案需由地方官、书吏、士绅、宗族、商人与劳动者共同转译；不同家庭面对的税役、婚姻、迁徙和安全风险并不相等。材料只让我们看见其中有限的记录者。承认可见与不可见的界线，才能使制度史、文化史和区域史互相校正。

在账本条目\texttt{undefined}中，1465（成化元年）的“荆襄山地流民聚众反抗，朝廷调兵镇压并酝酿安置方案”发生于明与荆襄流民。本条把背景说明为战乱、赋役逃避和山地垦殖吸引大量流动人口，地方官府难以将其纳入既有里甲和卫所秩序。，并以军事镇压之外，朝廷转向以招抚、编户和设治处理人口流动；原杰抚治成为这一转折的关键。概括可见影响。要理解它的社会后果，不能止于宫廷或战场：土地、赋役、司法、道路、性别化家庭劳动、城市市场、书院寺观和边海网络都会影响地方如何接收或改写制度。

本条依据《宪宗实录》成化元年荆襄条；《明史·原杰传》；郧阳、襄阳地方志。这些来源能够较稳妥地支持所记年序、诏令、任命或特定地点的行动，但无法自动证明全国的财产、人口、识字、信仰或动机。账本保留的证据说明是“起事及朝廷处置可靠；“流民”“盗贼”是行政分类，参与者来源与规模不能由官军奏报直接推出。”。所以，官府标签、奏疏数字与后来的道德评语必须放回其文书目的，不能被写成不经分区核验的社会全貌。

这条事件更适合作为一组关系变化的锚点：中央方案需由地方官、书吏、士绅、宗族、商人与劳动者共同转译；不同家庭面对的税役、婚姻、迁徙和安全风险并不相等。材料只让我们看见其中有限的记录者。承认可见与不可见的界线，才能使制度史、文化史和区域史互相校正。

在账本条目\texttt{undefined}中，1466（成化二年）一月的“明军击败并俘获侯大狗，但叛乱仍在继续”发生于明与地方反抗、流动作战集团。本条把背景说明为前期边防、地方武装或跨区域资源调动的压力推动了该行动。，并以它改变了后续调兵、补给或地方秩序的条件；战果和伤亡须分开核对。概括可见影响。要理解它的社会后果，不能止于宫廷或战场：土地、赋役、司法、道路、性别化家庭劳动、城市市场、书院寺观和边海网络都会影响地方如何接收或改写制度。

本条依据《宪宗实录》成化二年条；《明史·兵志》及相关列传；边镇、地方志或对方材料。这些来源能够较稳妥地支持所记年序、诏令、任命或特定地点的行动，但无法自动证明全国的财产、人口、识字、信仰或动机。账本保留的证据说明是“译写候选以编年材料定位；实录记载反映朝廷选择，崇祯期须另以奏疏、地方及敌对政权材料交叉。”。所以，官府标签、奏疏数字与后来的道德评语必须放回其文书目的，不能被写成不经分区核验的社会全貌。

这条事件更适合作为一组关系变化的锚点：中央方案需由地方官、书吏、士绅、宗族、商人与劳动者共同转译；不同家庭面对的税役、婚姻、迁徙和安全风险并不相等。材料只让我们看见其中有限的记录者。承认可见与不可见的界线，才能使制度史、文化史和区域史互相校正。

在账本条目\texttt{undefined}中，1466（成化二年）的“明军斩建州左卫东山”发生于明。本条把背景说明为继承、官僚协调或皇权运作的压力使问题进入朝廷程序。，并以它影响后续人事与政策议程；动机和派系标签不能由单一叙事裁定。概括可见影响。要理解它的社会后果，不能止于宫廷或战场：土地、赋役、司法、道路、性别化家庭劳动、城市市场、书院寺观和边海网络都会影响地方如何接收或改写制度。

本条依据《宪宗实录》成化二年条；《明史·本纪》及相关列传；诏令、奏疏或会典版本。这些来源能够较稳妥地支持所记年序、诏令、任命或特定地点的行动，但无法自动证明全国的财产、人口、识字、信仰或动机。账本保留的证据说明是“译写候选以编年材料定位；实录记载反映朝廷选择，崇祯期须另以奏疏、地方及敌对政权材料交叉。”。所以，官府标签、奏疏数字与后来的道德评语必须放回其文书目的，不能被写成不经分区核验的社会全貌。

这条事件更适合作为一组关系变化的锚点：中央方案需由地方官、书吏、士绅、宗族、商人与劳动者共同转译；不同家庭面对的税役、婚姻、迁徙和安全风险并不相等。材料只让我们看见其中有限的记录者。承认可见与不可见的界线，才能使制度史、文化史和区域史互相校正。

在账本条目\texttt{undefined}中，1466（成化二年）的“湖南及川黔边境苗族叛乱遭镇压”发生于明与地方反抗、流动作战集团。本条把背景说明为继承、官僚协调或皇权运作的压力使问题进入朝廷程序。，并以它影响后续人事与政策议程；动机和派系标签不能由单一叙事裁定。概括可见影响。要理解它的社会后果，不能止于宫廷或战场：土地、赋役、司法、道路、性别化家庭劳动、城市市场、书院寺观和边海网络都会影响地方如何接收或改写制度。

本条依据《宪宗实录》成化二年条；《明史·本纪》及相关列传；诏令、奏疏或会典版本。这些来源能够较稳妥地支持所记年序、诏令、任命或特定地点的行动，但无法自动证明全国的财产、人口、识字、信仰或动机。账本保留的证据说明是“译写候选以编年材料定位；实录记载反映朝廷选择，崇祯期须另以奏疏、地方及敌对政权材料交叉。”。所以，官府标签、奏疏数字与后来的道德评语必须放回其文书目的，不能被写成不经分区核验的社会全貌。

这条事件更适合作为一组关系变化的锚点：中央方案需由地方官、书吏、士绅、宗族、商人与劳动者共同转译；不同家庭面对的税役、婚姻、迁徙和安全风险并不相等。材料只让我们看见其中有限的记录者。承认可见与不可见的界线，才能使制度史、文化史和区域史互相校正。

在账本条目\texttt{undefined}中，1466（成化二年）的“刘通在襄阳附近起兵，兵败”发生于明与地方反抗、流动作战集团。本条把背景说明为前期边防、地方武装或跨区域资源调动的压力推动了该行动。，并以它改变了后续调兵、补给或地方秩序的条件；战果和伤亡须分开核对。概括可见影响。要理解它的社会后果，不能止于宫廷或战场：土地、赋役、司法、道路、性别化家庭劳动、城市市场、书院寺观和边海网络都会影响地方如何接收或改写制度。

本条依据《宪宗实录》成化二年条；《明史·兵志》及相关列传；边镇、地方志或对方材料。这些来源能够较稳妥地支持所记年序、诏令、任命或特定地点的行动，但无法自动证明全国的财产、人口、识字、信仰或动机。账本保留的证据说明是“译写候选以编年材料定位；实录记载反映朝廷选择，崇祯期须另以奏疏、地方及敌对政权材料交叉。”。所以，官府标签、奏疏数字与后来的道德评语必须放回其文书目的，不能被写成不经分区核验的社会全貌。

这条事件更适合作为一组关系变化的锚点：中央方案需由地方官、书吏、士绅、宗族、商人与劳动者共同转译；不同家庭面对的税役、婚姻、迁徙和安全风险并不相等。材料只让我们看见其中有限的记录者。承认可见与不可见的界线，才能使制度史、文化史和区域史互相校正。

在账本条目\texttt{undefined}中，1466（成化二年）的“原杰奉命抚治荆襄，招抚、编置部分流民并核定屯种”发生于明与荆襄流民。本条把背景说明为前一年军事处置未能解决流民生计和地方治理的持续问题。，并以招抚与编置使国家控制由临时征剿转向常设行政，为郧阳府设立和区域赋役重整铺路。概括可见影响。要理解它的社会后果，不能止于宫廷或战场：土地、赋役、司法、道路、性别化家庭劳动、城市市场、书院寺观和边海网络都会影响地方如何接收或改写制度。

本条依据《宪宗实录》成化二年条；《明史·原杰传》；《郧阳府志》有关设治和流民记载。这些来源能够较稳妥地支持所记年序、诏令、任命或特定地点的行动，但无法自动证明全国的财产、人口、识字、信仰或动机。账本保留的证据说明是“原杰奉命与安置措施可证；所报户口、田亩和安置成效是行政统计，不能等同全部流民的实际经历。”。所以，官府标签、奏疏数字与后来的道德评语必须放回其文书目的，不能被写成不经分区核验的社会全貌。

这条事件更适合作为一组关系变化的锚点：中央方案需由地方官、书吏、士绅、宗族、商人与劳动者共同转译；不同家庭面对的税役、婚姻、迁徙和安全风险并不相等。材料只让我们看见其中有限的记录者。承认可见与不可见的界线，才能使制度史、文化史和区域史互相校正。

在账本条目\texttt{undefined}中，1467（成化三年）的“明朝鲜联军击败建州女真并杀死李曼珠”发生于明与海上、朝贡相关政权。本条把背景说明为朝贡名义、边境安全与港口贸易的利益使交涉持续发生。，并以它为后续册封、互市或海防安排留下文书与跨政权比较线索。概括可见影响。要理解它的社会后果，不能止于宫廷或战场：土地、赋役、司法、道路、性别化家庭劳动、城市市场、书院寺观和边海网络都会影响地方如何接收或改写制度。

本条依据《宪宗实录》成化三年条；《明史·外国传》；《朝鲜王朝实录》、琉球《历代宝案》或海外档案。这些来源能够较稳妥地支持所记年序、诏令、任命或特定地点的行动，但无法自动证明全国的财产、人口、识字、信仰或动机。账本保留的证据说明是“译写候选以编年材料定位；实录记载反映朝廷选择，崇祯期须另以奏疏、地方及敌对政权材料交叉。”。所以，官府标签、奏疏数字与后来的道德评语必须放回其文书目的，不能被写成不经分区核验的社会全貌。

这条事件更适合作为一组关系变化的锚点：中央方案需由地方官、书吏、士绅、宗族、商人与劳动者共同转译；不同家庭面对的税役、婚姻、迁徙和安全风险并不相等。材料只让我们看见其中有限的记录者。承认可见与不可见的界线，才能使制度史、文化史和区域史互相校正。

\subsection{法律、家庭与社会关系}

在账本条目\texttt{undefined}中，1467（成化三年）的“科举、学校和法律条文在本年制度材料中可追踪，不能据规范文本推定州县实践一致。”发生于明。本条把背景说明为制度、教育或知识传播的需要推动了相关行动或文本形成。，并以它提供制度和传播史的锚点，不能据此推断社会整体接受。概括可见影响。要理解它的社会后果，不能止于宫廷或战场：土地、赋役、司法、道路、性别化家庭劳动、城市市场、书院寺观和边海网络都会影响地方如何接收或改写制度。

本条依据《宪宗实录》成化三年条；《明史·选举志》或《艺文志》；文集、碑刻或书目材料。这些来源能够较稳妥地支持所记年序、诏令、任命或特定地点的行动，但无法自动证明全国的财产、人口、识字、信仰或动机。账本保留的证据说明是“桥接条目只标示可回查的年度问题与材料链；实录有中央选择性，崇祯期尤其需要奏疏、地方、朝鲜及满文材料交叉。”。所以，官府标签、奏疏数字与后来的道德评语必须放回其文书目的，不能被写成不经分区核验的社会全貌。

这条事件更适合作为一组关系变化的锚点：中央方案需由地方官、书吏、士绅、宗族、商人与劳动者共同转译；不同家庭面对的税役、婚姻、迁徙和安全风险并不相等。材料只让我们看见其中有限的记录者。承认可见与不可见的界线，才能使制度史、文化史和区域史互相校正。

在账本条目\texttt{undefined}中，1467（成化三年）的“土木危机后的边镇整顿、京师防务或复辟前后军政调度在本年材料中构成连续问题。”发生于明。本条把背景说明为前期边防、地方武装或跨区域资源调动的压力推动了该行动。，并以它改变了后续调兵、补给或地方秩序的条件；战果和伤亡须分开核对。概括可见影响。要理解它的社会后果，不能止于宫廷或战场：土地、赋役、司法、道路、性别化家庭劳动、城市市场、书院寺观和边海网络都会影响地方如何接收或改写制度。

本条依据《宪宗实录》成化三年条；《明史·兵志》及相关列传；边镇、地方志或对方材料。这些来源能够较稳妥地支持所记年序、诏令、任命或特定地点的行动，但无法自动证明全国的财产、人口、识字、信仰或动机。账本保留的证据说明是“桥接条目只标示可回查的年度问题与材料链；实录有中央选择性，崇祯期尤其需要奏疏、地方、朝鲜及满文材料交叉。”。所以，官府标签、奏疏数字与后来的道德评语必须放回其文书目的，不能被写成不经分区核验的社会全貌。

这条事件更适合作为一组关系变化的锚点：中央方案需由地方官、书吏、士绅、宗族、商人与劳动者共同转译；不同家庭面对的税役、婚姻、迁徙和安全风险并不相等。材料只让我们看见其中有限的记录者。承认可见与不可见的界线，才能使制度史、文化史和区域史互相校正。

在账本条目\texttt{undefined}中，1467（成化三年）的“明军会同朝鲜军发动对建州女真据点的征讨，后世常称“成化犁庭””发生于明与建州女真诸部。本条把背景说明为建州诸部与辽东边市、朝贡及边境冲突相互交织，明朝和朝鲜均试图压缩其军事活动空间。，并以部分聚落和首领网络遭受重创，但建州诸部并未消失；辽东边政仍依赖招抚、贸易与军事威慑并用。概括可见影响。要理解它的社会后果，不能止于宫廷或战场：土地、赋役、司法、道路、性别化家庭劳动、城市市场、书院寺观和边海网络都会影响地方如何接收或改写制度。

本条依据《宪宗实录》成化三年建州条；《朝鲜王朝实录》世祖三年北征条；《明史·女真传》。这些来源能够较稳妥地支持所记年序、诏令、任命或特定地点的行动，但无法自动证明全国的财产、人口、识字、信仰或动机。账本保留的证据说明是“联合征讨发生可靠；“犁庭”、俘斩数字及其是否造成全面毁灭，均带有明朝与朝鲜军功叙事的夸张。”。所以，官府标签、奏疏数字与后来的道德评语必须放回其文书目的，不能被写成不经分区核验的社会全貌。

这条事件更适合作为一组关系变化的锚点：中央方案需由地方官、书吏、士绅、宗族、商人与劳动者共同转译；不同家庭面对的税役、婚姻、迁徙和安全风险并不相等。材料只让我们看见其中有限的记录者。承认可见与不可见的界线，才能使制度史、文化史和区域史互相校正。

在账本条目\texttt{undefined}中，1468（成化四年）五月的“蒙古人在固原起义”发生于明与蒙古诸部。本条把背景说明为继承、官僚协调或皇权运作的压力使问题进入朝廷程序。，并以它影响后续人事与政策议程；动机和派系标签不能由单一叙事裁定。概括可见影响。要理解它的社会后果，不能止于宫廷或战场：土地、赋役、司法、道路、性别化家庭劳动、城市市场、书院寺观和边海网络都会影响地方如何接收或改写制度。

本条依据《宪宗实录》成化四年条；《明史·本纪》及相关列传；诏令、奏疏或会典版本。这些来源能够较稳妥地支持所记年序、诏令、任命或特定地点的行动，但无法自动证明全国的财产、人口、识字、信仰或动机。账本保留的证据说明是“译写候选以编年材料定位；实录记载反映朝廷选择，崇祯期须另以奏疏、地方及敌对政权材料交叉。”。所以，官府标签、奏疏数字与后来的道德评语必须放回其文书目的，不能被写成不经分区核验的社会全貌。

这条事件更适合作为一组关系变化的锚点：中央方案需由地方官、书吏、士绅、宗族、商人与劳动者共同转译；不同家庭面对的税役、婚姻、迁徙和安全风险并不相等。材料只让我们看见其中有限的记录者。承认可见与不可见的界线，才能使制度史、文化史和区域史互相校正。

在账本条目\texttt{undefined}中，1468（成化四年）五月12日的“越南千军占领广西凭祥边境小镇”发生于明。本条把背景说明为朝贡名义、边境安全与港口贸易的利益使交涉持续发生。，并以它为后续册封、互市或海防安排留下文书与跨政权比较线索。概括可见影响。要理解它的社会后果，不能止于宫廷或战场：土地、赋役、司法、道路、性别化家庭劳动、城市市场、书院寺观和边海网络都会影响地方如何接收或改写制度。

本条依据《宪宗实录》成化四年条；《明史·外国传》；《朝鲜王朝实录》、琉球《历代宝案》或海外档案。这些来源能够较稳妥地支持所记年序、诏令、任命或特定地点的行动，但无法自动证明全国的财产、人口、识字、信仰或动机。账本保留的证据说明是“译写候选以编年材料定位；实录记载反映朝廷选择，崇祯期须另以奏疏、地方及敌对政权材料交叉。”。所以，官府标签、奏疏数字与后来的道德评语必须放回其文书目的，不能被写成不经分区核验的社会全貌。

这条事件更适合作为一组关系变化的锚点：中央方案需由地方官、书吏、士绅、宗族、商人与劳动者共同转译；不同家庭面对的税役、婚姻、迁徙和安全风险并不相等。材料只让我们看见其中有限的记录者。承认可见与不可见的界线，才能使制度史、文化史和区域史互相校正。

在账本条目\texttt{undefined}中，1468（成化四年）的“地方灾异、赈济与水利条目为本年社会史提供入口，地方志的修纂层次需要说明。”发生于明。本条把背景说明为自然风险与地方报告促成朝廷或地方的应对记录。，并以赈济、迁徙和损失的实际范围仍须以受灾地区材料分别核验。概括可见影响。要理解它的社会后果，不能止于宫廷或战场：土地、赋役、司法、道路、性别化家庭劳动、城市市场、书院寺观和边海网络都会影响地方如何接收或改写制度。

本条依据《宪宗实录》成化四年条；《明史·五行志》；受灾府县地方志与奏疏。这些来源能够较稳妥地支持所记年序、诏令、任命或特定地点的行动，但无法自动证明全国的财产、人口、识字、信仰或动机。账本保留的证据说明是“桥接条目只标示可回查的年度问题与材料链；实录有中央选择性，崇祯期尤其需要奏疏、地方、朝鲜及满文材料交叉。”。所以，官府标签、奏疏数字与后来的道德评语必须放回其文书目的，不能被写成不经分区核验的社会全貌。

这条事件更适合作为一组关系变化的锚点：中央方案需由地方官、书吏、士绅、宗族、商人与劳动者共同转译；不同家庭面对的税役、婚姻、迁徙和安全风险并不相等。材料只让我们看见其中有限的记录者。承认可见与不可见的界线，才能使制度史、文化史和区域史互相校正。

在账本条目\texttt{undefined}中，1468（成化四年）的“景泰、天顺时期的继承、人事与诏令链条可在本年材料中定位，政变责任不可由单一史书裁断。”发生于明。本条把背景说明为继承、官僚协调或皇权运作的压力使问题进入朝廷程序。，并以它影响后续人事与政策议程；动机和派系标签不能由单一叙事裁定。概括可见影响。要理解它的社会后果，不能止于宫廷或战场：土地、赋役、司法、道路、性别化家庭劳动、城市市场、书院寺观和边海网络都会影响地方如何接收或改写制度。

本条依据《宪宗实录》成化四年条；《明史·本纪》及相关列传；诏令、奏疏或会典版本。这些来源能够较稳妥地支持所记年序、诏令、任命或特定地点的行动，但无法自动证明全国的财产、人口、识字、信仰或动机。账本保留的证据说明是“桥接条目只标示可回查的年度问题与材料链；实录有中央选择性，崇祯期尤其需要奏疏、地方、朝鲜及满文材料交叉。”。所以，官府标签、奏疏数字与后来的道德评语必须放回其文书目的，不能被写成不经分区核验的社会全貌。

这条事件更适合作为一组关系变化的锚点：中央方案需由地方官、书吏、士绅、宗族、商人与劳动者共同转译；不同家庭面对的税役、婚姻、迁徙和安全风险并不相等。材料只让我们看见其中有限的记录者。承认可见与不可见的界线，才能使制度史、文化史和区域史互相校正。

在账本条目\texttt{undefined}中，1469（成化五年）的“固原蒙古叛乱被镇压”发生于明与蒙古诸部。本条把背景说明为前期边防、地方武装或跨区域资源调动的压力推动了该行动。，并以它改变了后续调兵、补给或地方秩序的条件；战果和伤亡须分开核对。概括可见影响。要理解它的社会后果，不能止于宫廷或战场：土地、赋役、司法、道路、性别化家庭劳动、城市市场、书院寺观和边海网络都会影响地方如何接收或改写制度。

本条依据《宪宗实录》成化五年条；《明史·兵志》及相关列传；边镇、地方志或对方材料。这些来源能够较稳妥地支持所记年序、诏令、任命或特定地点的行动，但无法自动证明全国的财产、人口、识字、信仰或动机。账本保留的证据说明是“译写候选以编年材料定位；实录记载反映朝廷选择，崇祯期须另以奏疏、地方及敌对政权材料交叉。”。所以，官府标签、奏疏数字与后来的道德评语必须放回其文书目的，不能被写成不经分区核验的社会全貌。

这条事件更适合作为一组关系变化的锚点：中央方案需由地方官、书吏、士绅、宗族、商人与劳动者共同转译；不同家庭面对的税役、婚姻、迁徙和安全风险并不相等。材料只让我们看见其中有限的记录者。承认可见与不可见的界线，才能使制度史、文化史和区域史互相校正。

在账本条目\texttt{undefined}中，1469（成化五年）的“景泰、天顺时期的继承、人事与诏令链条可在本年材料中定位，政变责任不可由单一史书裁断。”发生于明。本条把背景说明为继承、官僚协调或皇权运作的压力使问题进入朝廷程序。，并以它影响后续人事与政策议程；动机和派系标签不能由单一叙事裁定。概括可见影响。要理解它的社会后果，不能止于宫廷或战场：土地、赋役、司法、道路、性别化家庭劳动、城市市场、书院寺观和边海网络都会影响地方如何接收或改写制度。

本条依据《宪宗实录》成化五年条；《明史·本纪》及相关列传；诏令、奏疏或会典版本。这些来源能够较稳妥地支持所记年序、诏令、任命或特定地点的行动，但无法自动证明全国的财产、人口、识字、信仰或动机。账本保留的证据说明是“桥接条目只标示可回查的年度问题与材料链；实录有中央选择性，崇祯期尤其需要奏疏、地方、朝鲜及满文材料交叉。”。所以，官府标签、奏疏数字与后来的道德评语必须放回其文书目的，不能被写成不经分区核验的社会全貌。

这条事件更适合作为一组关系变化的锚点：中央方案需由地方官、书吏、士绅、宗族、商人与劳动者共同转译；不同家庭面对的税役、婚姻、迁徙和安全风险并不相等。材料只让我们看见其中有限的记录者。承认可见与不可见的界线，才能使制度史、文化史和区域史互相校正。

在账本条目\texttt{undefined}中，1469（成化五年）的“土木危机后的边镇整顿、京师防务或复辟前后军政调度在本年材料中构成连续问题。”发生于明。本条把背景说明为前期边防、地方武装或跨区域资源调动的压力推动了该行动。，并以它改变了后续调兵、补给或地方秩序的条件；战果和伤亡须分开核对。概括可见影响。要理解它的社会后果，不能止于宫廷或战场：土地、赋役、司法、道路、性别化家庭劳动、城市市场、书院寺观和边海网络都会影响地方如何接收或改写制度。

本条依据《宪宗实录》成化五年条；《明史·兵志》及相关列传；边镇、地方志或对方材料。这些来源能够较稳妥地支持所记年序、诏令、任命或特定地点的行动，但无法自动证明全国的财产、人口、识字、信仰或动机。账本保留的证据说明是“桥接条目只标示可回查的年度问题与材料链；实录有中央选择性，崇祯期尤其需要奏疏、地方、朝鲜及满文材料交叉。”。所以，官府标签、奏疏数字与后来的道德评语必须放回其文书目的，不能被写成不经分区核验的社会全貌。

这条事件更适合作为一组关系变化的锚点：中央方案需由地方官、书吏、士绅、宗族、商人与劳动者共同转译；不同家庭面对的税役、婚姻、迁徙和安全风险并不相等。材料只让我们看见其中有限的记录者。承认可见与不可见的界线，才能使制度史、文化史和区域史互相校正。

在账本条目\texttt{undefined}中，1469（成化五年）的“军饷、漕运和户籍赋役的本年记录揭示恢复秩序的行政成本，账面额与实收须分列。”发生于明。本条把背景说明为财政征解、交通工程或货币支付的压力使相关措施进入政务。，并以其法定安排不等于地方执行，后续须以地域材料追踪负担与成效。概括可见影响。要理解它的社会后果，不能止于宫廷或战场：土地、赋役、司法、道路、性别化家庭劳动、城市市场、书院寺观和边海网络都会影响地方如何接收或改写制度。

本条依据《宪宗实录》成化五年条；《明会典》相关条目；地方志、黄册/契约/河工材料。这些来源能够较稳妥地支持所记年序、诏令、任命或特定地点的行动，但无法自动证明全国的财产、人口、识字、信仰或动机。账本保留的证据说明是“桥接条目只标示可回查的年度问题与材料链；实录有中央选择性，崇祯期尤其需要奏疏、地方、朝鲜及满文材料交叉。”。所以，官府标签、奏疏数字与后来的道德评语必须放回其文书目的，不能被写成不经分区核验的社会全貌。

这条事件更适合作为一组关系变化的锚点：中央方案需由地方官、书吏、士绅、宗族、商人与劳动者共同转译；不同家庭面对的税役、婚姻、迁徙和安全风险并不相等。材料只让我们看见其中有限的记录者。承认可见与不可见的界线，才能使制度史、文化史和区域史互相校正。

\subsection{城市、文化与知识流通}

在账本条目\texttt{undefined}中，1470（成化六年）的“辽东巡抚陈越攻击女真人，向女真使馆索要贿赂”发生于明。本条把背景说明为继承、官僚协调或皇权运作的压力使问题进入朝廷程序。，并以它影响后续人事与政策议程；动机和派系标签不能由单一叙事裁定。概括可见影响。要理解它的社会后果，不能止于宫廷或战场：土地、赋役、司法、道路、性别化家庭劳动、城市市场、书院寺观和边海网络都会影响地方如何接收或改写制度。

本条依据《宪宗实录》成化六年条；《明史·本纪》及相关列传；诏令、奏疏或会典版本。这些来源能够较稳妥地支持所记年序、诏令、任命或特定地点的行动，但无法自动证明全国的财产、人口、识字、信仰或动机。账本保留的证据说明是“译写候选以编年材料定位；实录记载反映朝廷选择，崇祯期须另以奏疏、地方及敌对政权材料交叉。”。所以，官府标签、奏疏数字与后来的道德评语必须放回其文书目的，不能被写成不经分区核验的社会全貌。

这条事件更适合作为一组关系变化的锚点：中央方案需由地方官、书吏、士绅、宗族、商人与劳动者共同转译；不同家庭面对的税役、婚姻、迁徙和安全风险并不相等。材料只让我们看见其中有限的记录者。承认可见与不可见的界线，才能使制度史、文化史和区域史互相校正。

在账本条目\texttt{undefined}中，1470（成化六年）的“刘通叛军残部再次叛乱”发生于明与地方反抗、流动作战集团。本条把背景说明为继承、官僚协调或皇权运作的压力使问题进入朝廷程序。，并以它影响后续人事与政策议程；动机和派系标签不能由单一叙事裁定。概括可见影响。要理解它的社会后果，不能止于宫廷或战场：土地、赋役、司法、道路、性别化家庭劳动、城市市场、书院寺观和边海网络都会影响地方如何接收或改写制度。

本条依据《宪宗实录》成化六年条；《明史·本纪》及相关列传；诏令、奏疏或会典版本。这些来源能够较稳妥地支持所记年序、诏令、任命或特定地点的行动，但无法自动证明全国的财产、人口、识字、信仰或动机。账本保留的证据说明是“译写候选以编年材料定位；实录记载反映朝廷选择，崇祯期须另以奏疏、地方及敌对政权材料交叉。”。所以，官府标签、奏疏数字与后来的道德评语必须放回其文书目的，不能被写成不经分区核验的社会全貌。

这条事件更适合作为一组关系变化的锚点：中央方案需由地方官、书吏、士绅、宗族、商人与劳动者共同转译；不同家庭面对的税役、婚姻、迁徙和安全风险并不相等。材料只让我们看见其中有限的记录者。承认可见与不可见的界线，才能使制度史、文化史和区域史互相校正。

在账本条目\texttt{undefined}中，1470（成化六年）的“景泰、天顺时期的继承、人事与诏令链条可在本年材料中定位，政变责任不可由单一史书裁断。”发生于明。本条把背景说明为继承、官僚协调或皇权运作的压力使问题进入朝廷程序。，并以它影响后续人事与政策议程；动机和派系标签不能由单一叙事裁定。概括可见影响。要理解它的社会后果，不能止于宫廷或战场：土地、赋役、司法、道路、性别化家庭劳动、城市市场、书院寺观和边海网络都会影响地方如何接收或改写制度。

本条依据《宪宗实录》成化六年条；《明史·本纪》及相关列传；诏令、奏疏或会典版本。这些来源能够较稳妥地支持所记年序、诏令、任命或特定地点的行动，但无法自动证明全国的财产、人口、识字、信仰或动机。账本保留的证据说明是“桥接条目只标示可回查的年度问题与材料链；实录有中央选择性，崇祯期尤其需要奏疏、地方、朝鲜及满文材料交叉。”。所以，官府标签、奏疏数字与后来的道德评语必须放回其文书目的，不能被写成不经分区核验的社会全貌。

这条事件更适合作为一组关系变化的锚点：中央方案需由地方官、书吏、士绅、宗族、商人与劳动者共同转译；不同家庭面对的税役、婚姻、迁徙和安全风险并不相等。材料只让我们看见其中有限的记录者。承认可见与不可见的界线，才能使制度史、文化史和区域史互相校正。

在账本条目\texttt{undefined}中，1470（成化六年）的“蒙古诸部、朝鲜及西北绿洲的交涉在本年留下多方材料，应以具体使团和地名校正概称。”发生于明。本条把背景说明为朝贡名义、边境安全与港口贸易的利益使交涉持续发生。，并以它为后续册封、互市或海防安排留下文书与跨政权比较线索。概括可见影响。要理解它的社会后果，不能止于宫廷或战场：土地、赋役、司法、道路、性别化家庭劳动、城市市场、书院寺观和边海网络都会影响地方如何接收或改写制度。

本条依据《宪宗实录》成化六年条；《明史·外国传》；《朝鲜王朝实录》、琉球《历代宝案》或海外档案。这些来源能够较稳妥地支持所记年序、诏令、任命或特定地点的行动，但无法自动证明全国的财产、人口、识字、信仰或动机。账本保留的证据说明是“桥接条目只标示可回查的年度问题与材料链；实录有中央选择性，崇祯期尤其需要奏疏、地方、朝鲜及满文材料交叉。”。所以，官府标签、奏疏数字与后来的道德评语必须放回其文书目的，不能被写成不经分区核验的社会全貌。

这条事件更适合作为一组关系变化的锚点：中央方案需由地方官、书吏、士绅、宗族、商人与劳动者共同转译；不同家庭面对的税役、婚姻、迁徙和安全风险并不相等。材料只让我们看见其中有限的记录者。承认可见与不可见的界线，才能使制度史、文化史和区域史互相校正。

在账本条目\texttt{undefined}中，1471（成化七年）的“刘通的叛军被击败”发生于明与地方反抗、流动作战集团。本条把背景说明为前期边防、地方武装或跨区域资源调动的压力推动了该行动。，并以它改变了后续调兵、补给或地方秩序的条件；战果和伤亡须分开核对。概括可见影响。要理解它的社会后果，不能止于宫廷或战场：土地、赋役、司法、道路、性别化家庭劳动、城市市场、书院寺观和边海网络都会影响地方如何接收或改写制度。

本条依据《宪宗实录》成化七年条；《明史·兵志》及相关列传；边镇、地方志或对方材料。这些来源能够较稳妥地支持所记年序、诏令、任命或特定地点的行动，但无法自动证明全国的财产、人口、识字、信仰或动机。账本保留的证据说明是“译写候选以编年材料定位；实录记载反映朝廷选择，崇祯期须另以奏疏、地方及敌对政权材料交叉。”。所以，官府标签、奏疏数字与后来的道德评语必须放回其文书目的，不能被写成不经分区核验的社会全貌。

这条事件更适合作为一组关系变化的锚点：中央方案需由地方官、书吏、士绅、宗族、商人与劳动者共同转译；不同家庭面对的税役、婚姻、迁徙和安全风险并不相等。材料只让我们看见其中有限的记录者。承认可见与不可见的界线，才能使制度史、文化史和区域史互相校正。

在账本条目\texttt{undefined}中，1471（成化七年）的“军饷、漕运和户籍赋役的本年记录揭示恢复秩序的行政成本，账面额与实收须分列。（材料核查重点：诏令或奏议的形成与发受链）”发生于明。本条把背景说明为财政征解、交通工程或货币支付的压力使相关措施进入政务。，并以其法定安排不等于地方执行，后续须以地域材料追踪负担与成效。概括可见影响。要理解它的社会后果，不能止于宫廷或战场：土地、赋役、司法、道路、性别化家庭劳动、城市市场、书院寺观和边海网络都会影响地方如何接收或改写制度。

本条依据《宪宗实录》成化七年条；《明会典》相关条目；地方志、黄册/契约/河工材料。这些来源能够较稳妥地支持所记年序、诏令、任命或特定地点的行动，但无法自动证明全国的财产、人口、识字、信仰或动机。账本保留的证据说明是“桥接条目只标示可回查的年度问题与材料链；实录有中央选择性，崇祯期尤其需要奏疏、地方、朝鲜及满文材料交叉。；本条特别核对诏令或奏议的形成与发受链。”。所以，官府标签、奏疏数字与后来的道德评语必须放回其文书目的，不能被写成不经分区核验的社会全貌。

这条事件更适合作为一组关系变化的锚点：中央方案需由地方官、书吏、士绅、宗族、商人与劳动者共同转译；不同家庭面对的税役、婚姻、迁徙和安全风险并不相等。材料只让我们看见其中有限的记录者。承认可见与不可见的界线，才能使制度史、文化史和区域史互相校正。

在账本条目\texttt{undefined}中，1471（成化七年）的“军饷、漕运和户籍赋役的本年记录揭示恢复秩序的行政成本，账面额与实收须分列。（材料核查重点：报称信息的来源、抄传与行政分类）”发生于明。本条把背景说明为财政征解、交通工程或货币支付的压力使相关措施进入政务。，并以其法定安排不等于地方执行，后续须以地域材料追踪负担与成效。概括可见影响。要理解它的社会后果，不能止于宫廷或战场：土地、赋役、司法、道路、性别化家庭劳动、城市市场、书院寺观和边海网络都会影响地方如何接收或改写制度。

本条依据《宪宗实录》成化七年条；《明会典》相关条目；地方志、黄册/契约/河工材料。这些来源能够较稳妥地支持所记年序、诏令、任命或特定地点的行动，但无法自动证明全国的财产、人口、识字、信仰或动机。账本保留的证据说明是“桥接条目只标示可回查的年度问题与材料链；实录有中央选择性，崇祯期尤其需要奏疏、地方、朝鲜及满文材料交叉。；本条特别核对报称信息的来源、抄传与行政分类。”。所以，官府标签、奏疏数字与后来的道德评语必须放回其文书目的，不能被写成不经分区核验的社会全貌。

这条事件更适合作为一组关系变化的锚点：中央方案需由地方官、书吏、士绅、宗族、商人与劳动者共同转译；不同家庭面对的税役、婚姻、迁徙和安全风险并不相等。材料只让我们看见其中有限的记录者。承认可见与不可见的界线，才能使制度史、文化史和区域史互相校正。

在账本条目\texttt{undefined}中，1471（成化七年）的“科举、学校和法律条文在本年制度材料中可追踪，不能据规范文本推定州县实践一致。”发生于明。本条把背景说明为制度、教育或知识传播的需要推动了相关行动或文本形成。，并以它提供制度和传播史的锚点，不能据此推断社会整体接受。概括可见影响。要理解它的社会后果，不能止于宫廷或战场：土地、赋役、司法、道路、性别化家庭劳动、城市市场、书院寺观和边海网络都会影响地方如何接收或改写制度。

本条依据《宪宗实录》成化七年条；《明史·选举志》或《艺文志》；文集、碑刻或书目材料。这些来源能够较稳妥地支持所记年序、诏令、任命或特定地点的行动，但无法自动证明全国的财产、人口、识字、信仰或动机。账本保留的证据说明是“桥接条目只标示可回查的年度问题与材料链；实录有中央选择性，崇祯期尤其需要奏疏、地方、朝鲜及满文材料交叉。”。所以，官府标签、奏疏数字与后来的道德评语必须放回其文书目的，不能被写成不经分区核验的社会全貌。

这条事件更适合作为一组关系变化的锚点：中央方案需由地方官、书吏、士绅、宗族、商人与劳动者共同转译；不同家庭面对的税役、婚姻、迁徙和安全风险并不相等。材料只让我们看见其中有限的记录者。承认可见与不可见的界线，才能使制度史、文化史和区域史互相校正。

在账本条目\texttt{undefined}中，1472（成化八年）的“蒙古诸部、朝鲜及西北绿洲的交涉在本年留下多方材料，应以具体使团和地名校正概称。（材料核查重点：诏令或奏议的形成与发受链）”发生于明。本条把背景说明为朝贡名义、边境安全与港口贸易的利益使交涉持续发生。，并以它为后续册封、互市或海防安排留下文书与跨政权比较线索。概括可见影响。要理解它的社会后果，不能止于宫廷或战场：土地、赋役、司法、道路、性别化家庭劳动、城市市场、书院寺观和边海网络都会影响地方如何接收或改写制度。

本条依据《宪宗实录》成化八年条；《明史·外国传》；《朝鲜王朝实录》、琉球《历代宝案》或海外档案。这些来源能够较稳妥地支持所记年序、诏令、任命或特定地点的行动，但无法自动证明全国的财产、人口、识字、信仰或动机。账本保留的证据说明是“桥接条目只标示可回查的年度问题与材料链；实录有中央选择性，崇祯期尤其需要奏疏、地方、朝鲜及满文材料交叉。；本条特别核对诏令或奏议的形成与发受链。”。所以，官府标签、奏疏数字与后来的道德评语必须放回其文书目的，不能被写成不经分区核验的社会全貌。

这条事件更适合作为一组关系变化的锚点：中央方案需由地方官、书吏、士绅、宗族、商人与劳动者共同转译；不同家庭面对的税役、婚姻、迁徙和安全风险并不相等。材料只让我们看见其中有限的记录者。承认可见与不可见的界线，才能使制度史、文化史和区域史互相校正。

在账本条目\texttt{undefined}中，1472（成化八年）的“蒙古诸部、朝鲜及西北绿洲的交涉在本年留下多方材料，应以具体使团和地名校正概称。（材料核查重点：报称信息的来源、抄传与行政分类）”发生于明。本条把背景说明为朝贡名义、边境安全与港口贸易的利益使交涉持续发生。，并以它为后续册封、互市或海防安排留下文书与跨政权比较线索。概括可见影响。要理解它的社会后果，不能止于宫廷或战场：土地、赋役、司法、道路、性别化家庭劳动、城市市场、书院寺观和边海网络都会影响地方如何接收或改写制度。

本条依据《宪宗实录》成化八年条；《明史·外国传》；《朝鲜王朝实录》、琉球《历代宝案》或海外档案。这些来源能够较稳妥地支持所记年序、诏令、任命或特定地点的行动，但无法自动证明全国的财产、人口、识字、信仰或动机。账本保留的证据说明是“桥接条目只标示可回查的年度问题与材料链；实录有中央选择性，崇祯期尤其需要奏疏、地方、朝鲜及满文材料交叉。；本条特别核对报称信息的来源、抄传与行政分类。”。所以，官府标签、奏疏数字与后来的道德评语必须放回其文书目的，不能被写成不经分区核验的社会全貌。

这条事件更适合作为一组关系变化的锚点：中央方案需由地方官、书吏、士绅、宗族、商人与劳动者共同转译；不同家庭面对的税役、婚姻、迁徙和安全风险并不相等。材料只让我们看见其中有限的记录者。承认可见与不可见的界线，才能使制度史、文化史和区域史互相校正。

在账本条目\texttt{undefined}中，1472（成化八年）的“地方灾异、赈济与水利条目为本年社会史提供入口，地方志的修纂层次需要说明。”发生于明。本条把背景说明为自然风险与地方报告促成朝廷或地方的应对记录。，并以赈济、迁徙和损失的实际范围仍须以受灾地区材料分别核验。概括可见影响。要理解它的社会后果，不能止于宫廷或战场：土地、赋役、司法、道路、性别化家庭劳动、城市市场、书院寺观和边海网络都会影响地方如何接收或改写制度。

本条依据《宪宗实录》成化八年条；《明史·五行志》；受灾府县地方志与奏疏。这些来源能够较稳妥地支持所记年序、诏令、任命或特定地点的行动，但无法自动证明全国的财产、人口、识字、信仰或动机。账本保留的证据说明是“桥接条目只标示可回查的年度问题与材料链；实录有中央选择性，崇祯期尤其需要奏疏、地方、朝鲜及满文材料交叉。”。所以，官府标签、奏疏数字与后来的道德评语必须放回其文书目的，不能被写成不经分区核验的社会全貌。

这条事件更适合作为一组关系变化的锚点：中央方案需由地方官、书吏、士绅、宗族、商人与劳动者共同转译；不同家庭面对的税役、婚姻、迁徙和安全风险并不相等。材料只让我们看见其中有限的记录者。承认可见与不可见的界线，才能使制度史、文化史和区域史互相校正。

\subsection{环境、边防与海外连接}

在账本条目\texttt{undefined}中，1472（成化八年）的“吐鲁番阿黑麻势力进攻哈密，明朝在东天山的册封与朝贡网络受到冲击”发生于明、哈密与吐鲁番。本条把背景说明为哈密王族继承脆弱，而吐鲁番寻求控制绿洲、贸易和使团通道；明廷远距支援受嘉峪关外补给限制。，并以明廷须在册封、关市限制和边防成本之间重新取舍；1473年的继嗣交涉未消除吐鲁番压力。概括可见影响。要理解它的社会后果，不能止于宫廷或战场：土地、赋役、司法、道路、性别化家庭劳动、城市市场、书院寺观和边海网络都会影响地方如何接收或改写制度。

本条依据《宪宗实录》成化八年哈密条；《明史·西域传》；《明实录》所载哈密使臣文书。这些来源能够较稳妥地支持所记年序、诏令、任命或特定地点的行动，但无法自动证明全国的财产、人口、识字、信仰或动机。账本保留的证据说明是“进攻及明廷接获边报可靠；哈密城控制、当地王族选择和吐鲁番军队规模主要来自外部奏报，不能视为完整现场记录。”。所以，官府标签、奏疏数字与后来的道德评语必须放回其文书目的，不能被写成不经分区核验的社会全貌。

这条事件更适合作为一组关系变化的锚点：中央方案需由地方官、书吏、士绅、宗族、商人与劳动者共同转译；不同家庭面对的税役、婚姻、迁徙和安全风险并不相等。材料只让我们看见其中有限的记录者。承认可见与不可见的界线，才能使制度史、文化史和区域史互相校正。

在账本条目\texttt{undefined}中，1473（成化九年）的“明军与蒙古盟友联合进攻哈密，但在蒙古人放弃他们后撤退”发生于明与蒙古诸部。本条把背景说明为前期边防、地方武装或跨区域资源调动的压力推动了该行动。，并以它改变了后续调兵、补给或地方秩序的条件；战果和伤亡须分开核对。概括可见影响。要理解它的社会后果，不能止于宫廷或战场：土地、赋役、司法、道路、性别化家庭劳动、城市市场、书院寺观和边海网络都会影响地方如何接收或改写制度。

本条依据《宪宗实录》成化九年条；《明史·兵志》及相关列传；边镇、地方志或对方材料。这些来源能够较稳妥地支持所记年序、诏令、任命或特定地点的行动，但无法自动证明全国的财产、人口、识字、信仰或动机。账本保留的证据说明是“译写候选以编年材料定位；实录记载反映朝廷选择，崇祯期须另以奏疏、地方及敌对政权材料交叉。”。所以，官府标签、奏疏数字与后来的道德评语必须放回其文书目的，不能被写成不经分区核验的社会全貌。

这条事件更适合作为一组关系变化的锚点：中央方案需由地方官、书吏、士绅、宗族、商人与劳动者共同转译；不同家庭面对的税役、婚姻、迁徙和安全风险并不相等。材料只让我们看见其中有限的记录者。承认可见与不可见的界线，才能使制度史、文化史和区域史互相校正。

在账本条目\texttt{undefined}中，1473（成化九年）的“科举、学校和法律条文在本年制度材料中可追踪，不能据规范文本推定州县实践一致。”发生于明。本条把背景说明为制度、教育或知识传播的需要推动了相关行动或文本形成。，并以它提供制度和传播史的锚点，不能据此推断社会整体接受。概括可见影响。要理解它的社会后果，不能止于宫廷或战场：土地、赋役、司法、道路、性别化家庭劳动、城市市场、书院寺观和边海网络都会影响地方如何接收或改写制度。

本条依据《宪宗实录》成化九年条；《明史·选举志》或《艺文志》；文集、碑刻或书目材料。这些来源能够较稳妥地支持所记年序、诏令、任命或特定地点的行动，但无法自动证明全国的财产、人口、识字、信仰或动机。账本保留的证据说明是“桥接条目只标示可回查的年度问题与材料链；实录有中央选择性，崇祯期尤其需要奏疏、地方、朝鲜及满文材料交叉。”。所以，官府标签、奏疏数字与后来的道德评语必须放回其文书目的，不能被写成不经分区核验的社会全貌。

这条事件更适合作为一组关系变化的锚点：中央方案需由地方官、书吏、士绅、宗族、商人与劳动者共同转译；不同家庭面对的税役、婚姻、迁徙和安全风险并不相等。材料只让我们看见其中有限的记录者。承认可见与不可见的界线，才能使制度史、文化史和区域史互相校正。

在账本条目\texttt{undefined}中，1473（成化九年）的“土木危机后的边镇整顿、京师防务或复辟前后军政调度在本年材料中构成连续问题。”发生于明。本条把背景说明为前期边防、地方武装或跨区域资源调动的压力推动了该行动。，并以它改变了后续调兵、补给或地方秩序的条件；战果和伤亡须分开核对。概括可见影响。要理解它的社会后果，不能止于宫廷或战场：土地、赋役、司法、道路、性别化家庭劳动、城市市场、书院寺观和边海网络都会影响地方如何接收或改写制度。

本条依据《宪宗实录》成化九年条；《明史·兵志》及相关列传；边镇、地方志或对方材料。这些来源能够较稳妥地支持所记年序、诏令、任命或特定地点的行动，但无法自动证明全国的财产、人口、识字、信仰或动机。账本保留的证据说明是“桥接条目只标示可回查的年度问题与材料链；实录有中央选择性，崇祯期尤其需要奏疏、地方、朝鲜及满文材料交叉。”。所以，官府标签、奏疏数字与后来的道德评语必须放回其文书目的，不能被写成不经分区核验的社会全貌。

这条事件更适合作为一组关系变化的锚点：中央方案需由地方官、书吏、士绅、宗族、商人与劳动者共同转译；不同家庭面对的税役、婚姻、迁徙和安全风险并不相等。材料只让我们看见其中有限的记录者。承认可见与不可见的界线，才能使制度史、文化史和区域史互相校正。

在账本条目\texttt{undefined}中，1473（成化九年）的“哈密王位继承与朝贡使团问题引发明廷交涉，吐鲁番势力介入东天山政治”发生于明与哈密、吐鲁番。本条把背景说明为哈密既是明朝西域朝贡体系的节点，也受本地王族和吐鲁番力量竞争牵动。，并以明廷不得不在册封、赏赐、关市与边防之间权衡；哈密问题此后反复成为西北边政负担。概括可见影响。要理解它的社会后果，不能止于宫廷或战场：土地、赋役、司法、道路、性别化家庭劳动、城市市场、书院寺观和边海网络都会影响地方如何接收或改写制度。

本条依据《宪宗实录》成化九年哈密条；《明史·西域传》；《明实录》所载哈密使臣文书。这些来源能够较稳妥地支持所记年序、诏令、任命或特定地点的行动，但无法自动证明全国的财产、人口、识字、信仰或动机。账本保留的证据说明是“哈密政局和明廷交涉可见于实录；当地王位传承、部众立场和吐鲁番影响范围在中原材料中较难确证。”。所以，官府标签、奏疏数字与后来的道德评语必须放回其文书目的，不能被写成不经分区核验的社会全貌。

这条事件更适合作为一组关系变化的锚点：中央方案需由地方官、书吏、士绅、宗族、商人与劳动者共同转译；不同家庭面对的税役、婚姻、迁徙和安全风险并不相等。材料只让我们看见其中有限的记录者。承认可见与不可见的界线，才能使制度史、文化史和区域史互相校正。

在账本条目\texttt{undefined}中，1474（成化十年）的“于子军指挥重建和扩建长城，南封鄂尔多斯”发生于明。本条把背景说明为继承、官僚协调或皇权运作的压力使问题进入朝廷程序。，并以它影响后续人事与政策议程；动机和派系标签不能由单一叙事裁定。概括可见影响。要理解它的社会后果，不能止于宫廷或战场：土地、赋役、司法、道路、性别化家庭劳动、城市市场、书院寺观和边海网络都会影响地方如何接收或改写制度。

本条依据《宪宗实录》成化十年条；《明史·本纪》及相关列传；诏令、奏疏或会典版本。这些来源能够较稳妥地支持所记年序、诏令、任命或特定地点的行动，但无法自动证明全国的财产、人口、识字、信仰或动机。账本保留的证据说明是“译写候选以编年材料定位；实录记载反映朝廷选择，崇祯期须另以奏疏、地方及敌对政权材料交叉。”。所以，官府标签、奏疏数字与后来的道德评语必须放回其文书目的，不能被写成不经分区核验的社会全貌。

这条事件更适合作为一组关系变化的锚点：中央方案需由地方官、书吏、士绅、宗族、商人与劳动者共同转译；不同家庭面对的税役、婚姻、迁徙和安全风险并不相等。材料只让我们看见其中有限的记录者。承认可见与不可见的界线，才能使制度史、文化史和区域史互相校正。

在账本条目\texttt{undefined}中，1474（成化十年）的“科举、学校和法律条文在本年制度材料中可追踪，不能据规范文本推定州县实践一致。”发生于明。本条把背景说明为制度、教育或知识传播的需要推动了相关行动或文本形成。，并以它提供制度和传播史的锚点，不能据此推断社会整体接受。概括可见影响。要理解它的社会后果，不能止于宫廷或战场：土地、赋役、司法、道路、性别化家庭劳动、城市市场、书院寺观和边海网络都会影响地方如何接收或改写制度。

本条依据《宪宗实录》成化十年条；《明史·选举志》或《艺文志》；文集、碑刻或书目材料。这些来源能够较稳妥地支持所记年序、诏令、任命或特定地点的行动，但无法自动证明全国的财产、人口、识字、信仰或动机。账本保留的证据说明是“桥接条目只标示可回查的年度问题与材料链；实录有中央选择性，崇祯期尤其需要奏疏、地方、朝鲜及满文材料交叉。”。所以，官府标签、奏疏数字与后来的道德评语必须放回其文书目的，不能被写成不经分区核验的社会全貌。

这条事件更适合作为一组关系变化的锚点：中央方案需由地方官、书吏、士绅、宗族、商人与劳动者共同转译；不同家庭面对的税役、婚姻、迁徙和安全风险并不相等。材料只让我们看见其中有限的记录者。承认可见与不可见的界线，才能使制度史、文化史和区域史互相校正。

在账本条目\texttt{undefined}中，1474（成化十年）的“地方灾异、赈济与水利条目为本年社会史提供入口，地方志的修纂层次需要说明。”发生于明。本条把背景说明为自然风险与地方报告促成朝廷或地方的应对记录。，并以赈济、迁徙和损失的实际范围仍须以受灾地区材料分别核验。概括可见影响。要理解它的社会后果，不能止于宫廷或战场：土地、赋役、司法、道路、性别化家庭劳动、城市市场、书院寺观和边海网络都会影响地方如何接收或改写制度。

本条依据《宪宗实录》成化十年条；《明史·五行志》；受灾府县地方志与奏疏。这些来源能够较稳妥地支持所记年序、诏令、任命或特定地点的行动，但无法自动证明全国的财产、人口、识字、信仰或动机。账本保留的证据说明是“桥接条目只标示可回查的年度问题与材料链；实录有中央选择性，崇祯期尤其需要奏疏、地方、朝鲜及满文材料交叉。”。所以，官府标签、奏疏数字与后来的道德评语必须放回其文书目的，不能被写成不经分区核验的社会全貌。

这条事件更适合作为一组关系变化的锚点：中央方案需由地方官、书吏、士绅、宗族、商人与劳动者共同转译；不同家庭面对的税役、婚姻、迁徙和安全风险并不相等。材料只让我们看见其中有限的记录者。承认可见与不可见的界线，才能使制度史、文化史和区域史互相校正。

在账本条目\texttt{undefined}中，1474（成化十年）的“景泰、天顺时期的继承、人事与诏令链条可在本年材料中定位，政变责任不可由单一史书裁断。”发生于明。本条把背景说明为继承、官僚协调或皇权运作的压力使问题进入朝廷程序。，并以它影响后续人事与政策议程；动机和派系标签不能由单一叙事裁定。概括可见影响。要理解它的社会后果，不能止于宫廷或战场：土地、赋役、司法、道路、性别化家庭劳动、城市市场、书院寺观和边海网络都会影响地方如何接收或改写制度。

本条依据《宪宗实录》成化十年条；《明史·本纪》及相关列传；诏令、奏疏或会典版本。这些来源能够较稳妥地支持所记年序、诏令、任命或特定地点的行动，但无法自动证明全国的财产、人口、识字、信仰或动机。账本保留的证据说明是“桥接条目只标示可回查的年度问题与材料链；实录有中央选择性，崇祯期尤其需要奏疏、地方、朝鲜及满文材料交叉。”。所以，官府标签、奏疏数字与后来的道德评语必须放回其文书目的，不能被写成不经分区核验的社会全貌。

这条事件更适合作为一组关系变化的锚点：中央方案需由地方官、书吏、士绅、宗族、商人与劳动者共同转译；不同家庭面对的税役、婚姻、迁徙和安全风险并不相等。材料只让我们看见其中有限的记录者。承认可见与不可见的界线，才能使制度史、文化史和区域史互相校正。
